\documentclass{article}
\usepackage{amsmath,amssymb,amsthm}
\usepackage{algorithm,algorithmic}
\usepackage{enumitem}
\usepackage{hyperref}
\usepackage{bm}
\usepackage{color}
\newtheorem{theorem}{Theorem}
\newtheorem{lemma}{Lemma}
\newtheorem{assumption}{Assumption}
\newcommand{\E}{\mathbb{E}}
\newcommand{\R}{\mathbb{R}}

\begin{document}

\title{Randomized Coordinate Descent with Non‑Uniform Sampling\\
(Non‑convex, Coordinate‑wise Smooth Setting)}
\author{}
\date{}
\maketitle

%%%%%%%%%%%%%%%%%%%%%%%%%%%%%%%%%%%%%%%%%%%%%%%%%%%%%%%%%%%%%%%%%%%%%%%%%%%%
\section*{Algorithm Name}
\textbf{Randomized Coordinate Descent (RCD) with probabilities $p_i$}.  
At each iteration a single coordinate $i_k\in\{1,\dots,d\}$ is drawn from a
fixed distribution $\{p_i\}_{i=1}^d$ ($p_i>0$, $\sum_i p_i=1$).  
Only that coordinate is updated using the exact partial derivative and a
coordinate‑wise stepsize $1/L_i$ (or a generic stepsize $\eta_i$).

%%%%%%%%%%%%%%%%%%%%%%%%%%%%%%%%%%%%%%%%%%%%%%%%%%%%%%%%%%%%%%%%%%%%%%%%%%%%
\section*{Mathematical Formulation}
Let $f:\R^d\to\R$ be differentiable (possibly non‑convex).  
Denote the $i$‑th partial derivative by $\nabla_i f(w)$ and the $i$‑th
coordinate of $w$ by $w^{(i)}$.

\begin{enumerate}[label=\arabic*)]
\item \textbf{Sampling:} $i_k\sim\mathcal{P}$ with $\Pr(i_k=i)=p_i$, $p_i>0$.
\item \textbf{Update:}
\[
w_{k+1}^{(i)}=
\begin{cases}
w_k^{(i)} - \displaystyle\frac{1}{L_{i}}\,\nabla_i f(w_k), &\text{if }i=i_k,\\[2ex]
w_k^{(i)}, &\text{otherwise,}
\end{cases}
\tag{1}
\]
where $L_i>0$ is the coordinate‑wise smoothness constant (defined below).
\item \textbf{Other coordinates:} $w_{k+1}^{(j)}=w_k^{(j)}$ for $j\neq i_k$.
\item \textbf{Termination:} after $T$ iterations or when a stopping
criterion (e.g. $\| \nabla f(w_k)\| \le \varepsilon$) is satisfied.
\end{enumerate}

%%%%%%%%%%%%%%%%%%%%%%%%%%%%%%%%%%%%%%%%%%%%%%%%%%%%%%%%%%%%%%%%%%%%%%%%%%%%
\section*{Pseudocode}
\begin{algorithm}[H]
\caption{Randomized Coordinate Descent (RCD)}
\begin{algorithmic}[1]
\REQUIRE Initial point $w_0\in\R^d$, probabilities $\{p_i\}_{i=1}^d$, smoothness constants $\{L_i\}_{i=1}^d$, maximum iterations $T$.
\FOR{$k=0$ \TO $T-1$}
    \STATE Sample $i_k\in\{1,\dots,d\}$ with $\Pr(i_k=i)=p_i$.
    \STATE Compute partial gradient $g_i \gets \nabla_{i_k} f(w_k)$.
    \STATE Update coordinate:
    \[
    w_{k+1}^{(i_k)} \gets w_k^{(i_k)} - \frac{1}{L_{i_k}}\, g_i .
    \]
    \STATE For $j\neq i_k$, set $w_{k+1}^{(j)}\gets w_k^{(j)}$.
\ENDFOR
\RETURN $w_T$ (or the best iterate seen).
\end{algorithmic}
\end{algorithm}

%%%%%%%%%%%%%%%%%%%%%%%%%%%%%%%%%%%%%%%%%%%%%%%%%%%%%%%%%%%%%%%%%%%%%%%%%%%%
\section*{Dry Run Example}
Consider the one‑dimensional quadratic
\[
f(w) = (w-3)^2,\qquad w\in\R .
\]
Here $d=1$, $p_1=1$, and $L_1=2$ (since $\nabla^2 f(w)=2$).  
Take stepsize $1/L_1 = 0.5$ (equivalently $\eta=0.5$).  
Start from $w_0=0$ and run three iterations.

\begin{center}
\begin{tabular}{c|c|c|c}
$k$ & $g_k=\nabla f(w_k)$ & Update $w_{k+1}=w_k-\frac{1}{L_1}g_k$ & $f(w_{k+1})$\\ \hline
0 & $-6$ & $0-0.5(-6)=3$ & $0$\\
1 & $0$  & $3-0.5(0)=3$ & $0$\\
2 & $0$  & $3-0.5(0)=3$ & $0$
\end{tabular}
\end{center}

After the first iteration the algorithm lands exactly at the global
minimiser $w^\star=3$, and the objective stays at $0$ thereafter.

%%%%%%%%%%%%%%%%%%%%%%%%%%%%%%%%%%%%%%%%%%%%%%%%%%%%%%%%%%%%%%%%%%%%%%%%%%%%
\section*{Assumptions}
\begin{assumption}[Coordinate‑wise $L_i$‑smoothness]\label{ass:smooth}
For each coordinate $i\in\{1,\dots,d\}$ there exists $L_i>0$ such that
\[
\bigl| \nabla_i f(x+ \alpha e_i) - \nabla_i f(x) \bigr|
\;\le\; L_i |\alpha|
\qquad\forall x\in\R^d,\;\forall \alpha\in\R,
\]
where $e_i$ is the $i$‑th unit vector. Equivalently,
\begin{equation}
f(x+\alpha e_i) \le f(x) + \alpha \nabla_i f(x) + \frac{L_i}{2}\alpha^2 .
\tag{2}
\end{equation}
\end{assumption}

\begin{assumption}[Bounded variance of stochastic partial gradients]\label{ass:bounded}
The algorithm uses the exact partial derivative, hence
\[
\operatorname{Var}\bigl(\nabla_{i_k} f(w_k)\mid w_k\bigr)=0 .
\]
\end{assumption}

\begin{assumption}[Positive sampling probabilities]\label{ass:prob}
\[
p_i>0,\qquad \sum_{i=1}^d p_i =1 .
\]
\end{assumption}

No convexity is assumed; $f$ may be non‑convex.

%%%%%%%%%%%%%%%%%%%%%%%%%%%%%%%%%%%%%%%%%%%%%%%%%%%%%%%%%%%%%%%%%%%%%%%%%%%%
\section*{Main Convergence Theorem}
\begin{theorem}[Expected squared gradient norm decay]\label{thm:main}
Let $\{w_k\}_{k\ge0}$ be generated by the RCD algorithm under
Assumptions~\ref{ass:smooth}--\ref{ass:prob}. Define
\[
\Phi_k \;:=\; f(w_k)-f^\star,\qquad 
f^\star:=\inf_{w\in\R^d} f(w) .
\]
Then for any $T\ge1$,
\begin{equation}
\frac{1}{T}\sum_{k=0}^{T-1}\E\!\bigl[ \|\nabla f(w_k)\|^2 \bigr]
\;\le\;
\frac{2\,\bigl(\max_i L_i\bigr)}{\displaystyle \min_i p_i}\;
\frac{ \Phi_0 }{T}.
\tag{3}
\end{equation}
Consequently,
\[
\min_{0\le k<T}\E\!\bigl[ \|\nabla f(w_k)\|^2 \bigr]
\;=\; O\!\left(\frac{1}{T}\right).
\]
\end{theorem}

\begin{theorem}[Almost‑sure convergence to stationary points]\label{thm:as}
If $\sum_{k=0}^\infty \E[\|\nabla f(w_k)\|^2] <\infty$ (which follows from (3)),
then $\|\nabla f(w_k)\|\to0$ almost surely and every accumulation point of
$\{w_k\}$ is a first‑order stationary point of $f$.
\end{theorem}

%%%%%%%%%%%%%%%%%%%%%%%%%%%%%%%%%%%%%%%%%%%%%%%%%%%%%%%%%%%%%%%%%%%%%%%%%%%%
\section*{Theoretical Properties}
\begin{itemize}
\item \textbf{Stability w.r.t.\ sampling distribution.} The bound (3) contains
$\displaystyle \frac{1}{\min_i p_i}$, showing that very small probabilities
slow down convergence. Uniform sampling ($p_i=1/d$) yields a factor $d$.
\item \textbf{Hyper‑parameter sensitivity.} Only the smoothness constants $L_i$
appear; the algorithm is \emph{stepsize‑free} once $L_i$ are known. Over‑estimating
$L_i$ makes the method more conservative but still guarantees (3).
\item \textbf{Comparison to full gradient descent.} Full GD with stepsize $1/L_{\max}$
gives $O(1/T)$ decay of $\|\nabla f\|^2$ with a factor $L_{\max}$. RCD replaces $L_{\max}$
by $\max_i L_i$ and introduces the sampling factor $1/\min_i p_i$, but each iteration
is $1/d$ cheaper.
\end{itemize}

\section*{Discussion}
The theorem proves that, despite the lack of convexity, the expected squared norm
of the full gradient decays at the optimal $O(1/T)$ rate for first‑order
methods applied to smooth (possibly non‑convex) objectives. The proof crucially
uses the weighting by $1/p_i$ when taking expectations, which allows the
partial‑gradient contributions to be summed into the full gradient norm.

%%%%%%%%%%%%%%%%%%%%%%%%%%%%%%%%%%%%%%%%%%%%%%%%%%%%%%%%%%%%%%%%%%%%%%%%%%%%
\section*{Preliminaries (Definitions and Lemmas)}
\begin{lemma}[One‑step progress under coordinate smoothness]\label{lem:one-step}
For any $w\in\R^d$ and any coordinate $i$, the update
\[
w^{+}=w - \frac{1}{L_i}\nabla_i f(w) e_i
\]
satisfies
\[
f(w^{+}) \le f(w) - \frac{1}{2L_i}\bigl(\nabla_i f(w)\bigr)^2 .
\tag{4}
\]
\end{lemma}
\begin{proof}
Apply the smoothness inequality (2) with $\alpha=-\nabla_i f(w)/L_i$:
\[
\begin{aligned}
f\!\bigl(w-\tfrac{1}{L_i}\nabla_i f(w) e_i \bigr)
&\le f(w) -\frac{1}{L_i}\bigl(\nabla_i f(w)\bigr)^2
   + \frac{L_i}{2}\Bigl(\frac{\nabla_i f(w)}{L_i}\Bigr)^2    \\
&= f(w) - \frac{1}{2L_i}\bigl(\nabla_i f(w)\bigr)^2 .
\end{aligned}
\]
\end{proof}

%%%%%%%%%%%%%%%%%%%%%%%%%%%%%%%%%%%%%%%%%%%%%%%%%%%%%%%%%%%%%%%%%%%%%%%%%%%%
\section*{Main Proof (Step‑by‑Step Derivation)}
We prove Theorem~\ref{thm:main}. Every non‑trivial manipulation is annotated
with a line number ``[Step N]''.

\begin{proof}
Let $w_k$ be the iterate at time $k$ and $i_k$ the sampled coordinate.
From Lemma~\ref{lem:one-step} applied to $i=i_k$ we have
\begin{equation}
f(w_{k+1}) \le f(w_k) - 
\frac{1}{2L_{i_k}}\bigl(\nabla_{i_k} f(w_k)\bigr)^2 .
\tag{5}
\end{equation}
\textbf{[Step 1]} Take total expectation conditioned on $w_k$ and use the law of
total expectation:
\[
\begin{aligned}
\E\!\bigl[f(w_{k+1})\mid w_k\bigr]
&= \sum_{i=1}^d p_i\,
   \Bigl( f(w_k) - \frac{1}{2L_i}\bigl(\nabla_i f(w_k)\bigr)^2 \Bigr) \\
&= f(w_k) - \frac12\sum_{i=1}^d p_i\frac{1}{L_i}
   \bigl(\nabla_i f(w_k)\bigr)^2 .
\end{aligned}
\tag{6}
\]

\textbf{[Step 2]} Rearrange (6) to isolate the weighted squared gradient:
\[
\frac12\sum_{i=1}^d p_i\frac{1}{L_i}
   \bigl(\nabla_i f(w_k)\bigr)^2
   \;\le\; f(w_k)-\E\!\bigl[f(w_{k+1})\mid w_k\bigr] .
\tag{7}
\]

\textbf{[Step 3]} Apply the elementary inequality
\[
\frac{1}{p_i}\ge \frac{1}{\min_j p_j}\qquad\forall i,
\]
multiply both sides of (7) by $\displaystyle \frac{2\max_j L_j}{\min_j p_j}$,
and use $\max_j L_j\ge L_i$ for each $i$:
\[
\begin{aligned}
\sum_{i=1}^d \bigl(\nabla_i f(w_k)\bigr)^2
&\le
\frac{2\max_j L_j}{\min_j p_j}
\Bigl( f(w_k)-\E\!\bigl[f(w_{k+1})\mid w_k\bigr] \Bigr) .
\end{aligned}
\tag{8}
\]

\textbf{[Step 4}] Take full expectation over all randomness up to iteration $k$
(and denote $\E_k[\cdot]=\E[\cdot\mid w_k]$):
\[
\E\!\bigl[\|\nabla f(w_k)\|^2\bigr]
\;\le\;
\frac{2\max_j L_j}{\min_j p_j}
\Bigl( \E\!\bigl[f(w_k)\bigr]-\E\!\bigl[f(w_{k+1})\bigr] \Bigr) .
\tag{9}
\]

\textbf{[Step 5]} Sum inequality (9) over $k=0,\dots,T-1$:
\[
\sum_{k=0}^{T-1}\E\!\bigl[\|\nabla f(w_k)\|^2\bigr]
\;\le\;
\frac{2\max_j L_j}{\min_j p_j}
\Bigl( \E\!\bigl[f(w_0)\bigr]-\E\!\bigl[f(w_T)\bigr] \Bigr) .
\tag{10}
\]

Since $f(w_T)\ge f^\star$, we may replace $-\E[f(w_T)]$ by $-f^\star$ and obtain
\[
\sum_{k=0}^{T-1}\E\!\bigl[\|\nabla f(w_k)\|^2\bigr]
\;\le\;
\frac{2\max_j L_j}{\min_j p_j}
\bigl( f(w_0)-f^\star \bigr) .
\tag{11}
\]

\textbf{[Step 6]} Divide both sides of (11) by $T$ to get the average bound:
\[
\frac{1}{T}\sum_{k=0}^{T-1}\E\!\bigl[\|\nabla f(w_k)\|^2\bigr]
\;\le\;
\frac{2\max_j L_j}{\min_j p_j}\;
\frac{ f(w_0)-f^\star }{T} .
\tag{12}
\]

Equation (12) is exactly the claimed bound (3). ∎
\end{proof}

%%%%%%%%%%%%%%%%%%%%%%%%%%%%%%%%%%%%%%%%%%%%%%%%%%%%%%%%%%%%%%%%%%%%%%%%%%%%
\section*{Rate Derivation (With Constants)}
From (12) we read off the explicit constant
\[
C \;=\; \frac{2\max_i L_i}{\min_i p_i}\bigl( f(w_0)-f^\star \bigr) .
\]
Thus the algorithm enjoys an $O(1/T)$ decay of the expected squared
gradient norm, with the constant scaling linearly with the worst smoothness
parameter and inversely with the smallest sampling probability.

%%%%%%%%%%%%%%%%%%%%%%%%%%%%%%%%%%%%%%%%%%%%%%%%%%%%%%%%%%%%%%%%%%%%%%%%%%%%
\section*{Special Cases}
\begin{enumerate}[label=(\alph*)]
\item \textbf{Uniform sampling.} $p_i=1/d$ for all $i$.
   Then (3) becomes
   \[
   \frac{1}{T}\sum_{k=0}^{T-1}\E[\|\nabla f(w_k)\|^2]
   \le 2d\,\max_i L_i\,
   \frac{f(w_0)-f^\star}{T}.
   \]
   The factor $d$ reflects the fact that only one out of $d$ coordinates is
   updated per iteration.

\item \textbf{Importance sampling.} Choose $p_i\propto L_i$.
   Then $\frac{\max_i L_i}{\min_i p_i}= \frac{\max_i L_i}{\min_i (L_i/\sum_j L_j)}
   = \frac{\max_i L_i\,\sum_j L_j}{\min_i L_i}
   \le d\,\max_i L_i$.
   Hence the bound is never worse than uniform sampling and can be
   substantially better when the $L_i$'s are highly heterogeneous.

\item \textbf{Exact line‑search on the chosen coordinate.}
   If we replace the fixed stepsize $1/L_i$ by the exact minimiser along $e_i$,
   the one‑step inequality (4) becomes
   \[
   f(w_{k+1}) \le f(w_k) - \frac{1}{2L_i}\bigl(\nabla_i f(w_k)\bigr)^2,
   \]
   i.e. the same bound, so the convergence rate remains $O(1/T)$.

\end{enumerate}

%%%%%%%%%%%%%%%%%%%%%%%%%%%%%%%%%%%%%%%%%%%%%%%%%%%%%%%%%%%%%%%%%%%%%%%%%%%%
\section*{References}
\begin{enumerate}
\item Nesterov, Y. (2012). \emph{Efficiency of Coordinate Descent Methods on Huge-Scale Optimization Problems}. SIAM J. Optim.
\item Richtárik, P., & Takáč, M. (2014). \emph{Iteration Complexity of Random Coordinate Descent Methods for Composite Optimization}. Math. Program.
\end{enumerate}

\end{document}